\section{Fazit und Ausblick} \label{fazit}

Dieses Kapitel fasst die Ergebnisse zusammen und ordnet weitere Forschungsschritte ein.

\subsection{Zusammenfassung der Ergebnisse}

Die Projektarbeit verfolgt das Ziel, einen nachvollziehbaren Ansatz für KI-gestützte Code-Migration im Sinne des Software-Re-Engineering zu entwickeln. Im Mittelpunkt steht eine subagent-basierte Pipeline, die Quellcode analysiert, in einzelne Funktionalitäten zerlegt, daraus textuelle Spezifikationen erzeugt und darauf aufbauend Zielcode generiert. Der PoC zeigt für das Fallbeispiel, dass sich eine Migration über eine explizite Spezifikationsebene als Zwischenartefakt durchführen lässt und dass Abweichungen im Zielsystem auf einzelne Spezifikationen und Pipelineschritte zurückgeführt werden können. Ob die Spezifikationsebene die Risiken nicht-deterministischer LLM-Ausgaben tatsächlich reduziert, lässt sich ohne Vergleichsbedingung nicht belegen. Im manuellen I/O-Vergleich des zweiten Durchlaufs (Abschnitt~\ref{sec:io-befunde}) waren sieben von acht inhaltlichen Vergleichsfällen semantisch identisch. Zugleich wies er zwei methodisch eigenständige Defektklassen aus, die in zwei Skill-Iterationen aufgearbeitet wurden (Abschnitte~\ref{sec:builder-gap} und~\ref{sec:datafluss-gap}). Im dritten Durchlauf traten beide Defekte nicht mehr auf: 60 von 62 Aufrufen lieferten byteweise identische Ausgaben, die beiden Abweichungen betreffen denselben Fehlerfall (Abschnitt~\ref{sec:dritter-lauf}). Da sich in diesem Durchlauf neben den Skills auch Modell, OpenCode-Version und Code-Intelligence-Werkzeug geändert haben, lässt sich die Verbesserung nicht allein den Skill-Iterationen zuschreiben. Die Befunde beruhen auf einem einzigen Fallbeispiel und einem Lauf je Konfiguration (vgl.\ Abschnitt~\ref{sec:validitaet}). Mit OpenCode als Laufzeitumgebung und MCP-gebundenen Werkzeugen für Code-Intelligence und Dokumentationszugriff ist der PoC auf offenen Standards aufgebaut. Eine Übertragung in weitere Agentenumgebungen erscheint dadurch möglich, wurde aber nicht geprüft.

\subsection{Kritische Bewertung des spezifikationsgetriebenen Ansatzes}

Kritisch zu bewerten ist insbesondere der zusätzliche Aufwand, der durch die Erstellung und Prüfung der Spezifikation entsteht. Dem steht der konzeptionelle Vorteil gegenüber, dass Annahmen explizit gemacht, generierte Artefakte besser kontrolliert und Halluzinationsrisiken potenziell begrenzt werden können. Diese Vorteile sind konzeptionell plausibel und durch verwandte Arbeiten gestützt \parencite{midolo2026promptguidelines}. Der durchgeführte I/O-Vergleich konkretisiert sie für das vorliegende Migrationsszenario in zweierlei Hinsicht, wenn auch mit der in Abschnitt~\ref{sec:validitaet} beschriebenen begrenzten Reichweite: Einerseits produzierten die ersten Pipelinedurchläufe eine Klasse von Defekten, die artefaktinterne Tests und Builds nicht erkennen ließen und die erst durch beobachtungsbasiertes Differential Testing sichtbar wurden. Andererseits deutet er darauf hin, dass sich diese Defektklassen abstrakt benennen lassen (Verkabelungs- und Datenflusslücke). Sie sind in den Pipeline-Skills strukturell adressierbar, ohne dass die Agentenprofile domänenspezifische Sonderbehandlungen benötigen. Die in Kapitel~\ref{implementierung} formulierte Frage, ob der Rollenverbund mit Verifier ohne zusätzliche Evaluatoren-Rolle ausreicht, ist damit nur eingeschränkt zu bejahen: Solange die Verifier-Checkliste nur Artefakte und Builds prüft, bleiben Defekte im beobachtbaren Verhalten unentdeckt. Ein End-to-End-Vergleich gegen das Quellsystem gehört deshalb in die Prüfkette, sei es manuell wie im PoC oder automatisiert wie im Ausblick skizziert.

\subsection{Ausblick auf die Bachelorarbeit}

Für die Bachelorarbeit ist eine Erweiterung zum hybriden Ansatz aus Spec-Driven Development (SDD) und Model-Driven Development (MDD) vorgesehen. MDD bezeichnet dabei einen Ansatz, bei dem formale oder visuelle Modelle, etwa UML-Diagramme oder domänenspezifische Modellierungssprachen, als primäre Artefakte die Stelle textueller Beschreibungen einnehmen und direkt zur Code-Generierung genutzt werden \parencite{raibulet2017mdre}. Textuelle Spezifikationen sollen dann durch visuelle Modelle ergänzt werden, um strukturelle Beziehungen und Architekturaspekte noch präziser abbilden zu können. Darüber hinaus stellt sich die Frage, wie weit agentische Autonomie in der Codegenerierung sinnvoll ausgedehnt werden kann, ohne menschliche Kontroll- und Review-Schritte aufzugeben \parencite{boeckeler2025autonomy}.

Ein zusätzlicher Schwerpunkt der Bachelorarbeit soll in der Entwicklung eines reproduzierbaren Evaluationsrahmens liegen. Im Sinne eines Harness Engineering sollen Spezifikationsabgleich, Build-Ausführung, Testläufe und statische Analysen zu einer wiederholbaren Prüfkette zusammengeführt werden, um Halluzinationen, funktionale Abweichungen und Qualitätsverluste systematisch erfassen zu können. Methodisch knüpft dieser Rahmen an das klassische \emph{Differential Testing} an, bei dem zwei Implementierungen desselben Programms mit identischen Eingaben gegeneinander ausgeführt und ihre Ausgaben verglichen werden \parencite{mckeeman1998differentialtesting}. Als konkrete Operationalisierung dieses Prinzips für LLM-gestützte Code-Translation dient Fluorine \parencite{eniser2025realworld}, ein Cross-Language-Differenzfuzzer für automatisierte I/O-Äquivalenzprüfung zwischen Quell- und Zielimplementierung ohne handgeschriebene Testfälle. Dieser Ansatz zeigt, wie der im PoC manuell durchgeführte I/O-Vergleich in der Bachelorarbeit zu einer automatisierbaren Prüfkette ausgebaut werden könnte. Ein solcher Rahmen ist insbesondere für LLM-gestützte Softwareaufgaben relevant, weil die Güte der Ergebnisse nicht allein aus ihrer textuellen Plausibilität, sondern erst aus ihrer Ausführbarkeit und Konsistenz mit dem Quellkontext hervorgeht \parencite{yang2024sweagent,liu2024hallucode}. Zugleich argumentieren \textcite{korn2026reportingprompting}, dass Prompt- und Kontextgestaltung, Modellversionen und Threats to Validity explizit dokumentiert werden sollten, um Reproduzierbarkeit methodisch abzusichern.

\subsubsection{Skalierung des Wissenszugriffs}

Ein zentrales Forschungsthema ist die Skalierung des Wissenszugriffs bei großen Softwaresystemen. Dabei ist zu untersuchen, wie MCP-basierte Abrufmechanismen und strukturierte Wissensrepräsentationen auch bei umfangreichen Projektlandschaften effizient eingesetzt werden können. Neuere Vorschläge für Benchmark-Infrastrukturen im Software Engineering deuten zudem darauf hin, dass künftig stärker mit umfangreicherem Projektkontext und mehrdimensionalen Evaluationsmetriken gearbeitet werden sollte \parencite{rodriguezcardenas2026behelm}.

\subsubsection{Weitere Forschungsthemen}

Darüber hinaus stellt sich die Frage, wie die idiomatische Qualität des generierten Zielcodes systematisch bewertet werden kann. Rein funktionale Korrektheit reicht als Kriterium nicht aus, wenn Code zwar lauffähig, aber stilistisch weit vom Zielökosystem entfernt ist. Hierzu sollen geeignete Metriken entwickelt oder aus verwandten Arbeiten adaptiert werden. Schließlich eröffnet die Integration visueller Modelle zusätzliche Möglichkeiten für die strukturelle Analyse und Dokumentation von Migrationsentscheidungen, die im Rahmen des hybriden SDD-MDD-Ansatzes der Bachelorarbeit untersucht werden sollen.
